\chapter{การดีบั๊กขั้นสูงด้วย Chrome DevTools และการปรับแต่ง Core Web Vitals}

การส่งมอบเว็บแอปพลิเคชันสู่สายตาผู้ใช้ จำเป็นต้องผ่านการตรวจสอบประสิทธิภาพเชิงลึกตามมาตรฐานของ Google Core Web Vitals และการตรวจสอบความถูกต้องของการเข้าถึง

\section{การตรวจจับการรั่วไหลของหน่วยความจำใน JavaScript}

การเกิดหน่วยความจำรั่วไหล (Memory Leak) มักเกิดจาก 3 สาเหตุหลัก:
\begin{enumerate}
    \item \textbf{Detached DOM Nodes:} องค์ประกอบที่ถูกลบออกจากหน้าจอแต่ยังมีตัวแปรอ้างอิงอยู่
    \item \textbf{Uncleared Event Listeners:} ตัวดักจับอีเวนต์ที่ไม่ได้ทำการถอดออกเมื่อปิดคอมโพเนนต์
    \item \textbf{Retained Closures:} ฟังก์ชันปิดผนึกที่เก็บข้อมูลขนาดใหญ่ไว้ในขอบเขตการทำงาน
\end{enumerate}

การใช้เครื่องมือ Heap Profiler ใน Chrome DevTools จะช่วยแสดงกราฟลำดับการอ้างอิง (Retainer Tree) ทำให้ชี้เป้าวัตถุต้นเหตุได้อย่างแม่นยำ

\section{การปรับแต่งตัวชี้วัดความเร็ว Largest Contentful Paint}

เกณฑ์ของ Core Web Vitals กำหนดว่าค่า LCP ต้องน้อยกว่า $2.5\,\text{วินาที}$:

\begin{equation}
\text{LCP} = T_{\text{TTFB}} + T_{\text{LoadDelay}} + T_{\text{LoadTime}} + T_{\text{RenderDelay}} < 2.5\,\text{s}
\end{equation}

การเพิ่มคุณสมบัติ \texttt{fetchpriority="high"} ให้กับภาพฮีโร่และฟอนต์หลัก ร่วมกับการใช้เทคนิค Resource Preloading ช่วยลดค่า LCP ได้มากกว่าร้อยละ 40
